% !PRINTED 509
\weberpage{509}
%<<CONT
$Q$ ungeändert bleibt; und \gesperrt{zweitens} ist jede Function, die
durch $Q$ ungeändert bleibt, nach dem \gesperrt{Lagrange}'schen Satze in
$\Omega'$ enthalten. Dies sind aber [nach §. 149, a), b)] die charakteristischen
 Merkmale der \gesperrt{Galois}'schen Gruppe im Körper $\Omega'$.
Um also die \gesperrt{Galois}'sche Gruppe vom Grade $p$ auf den Grad $q$
zu reduciren, muss \gesperrt{eine} Wurzel einer Hülfsgleichung $j^{\text{ten}}$ Grades
in $\Omega$ adjungirt werden.

Den Satz 3. können wir auch so ausdrücken:

\gesperrt{Der Normalkörper $N = \Omega(\alpha, \alpha_1 \ldots \alpha_{m-1})$ ist ein Körper
$p^{\text{ten}}$ Grades über $\Omega$ und $q^{\text{ten}}$ Grades über $\Omega' = \Omega(\psi)$.}

Der Erniedrigung der Gruppe durch Adjunction von $\psi$ entspricht,
wie schon aus den allgemeinen Grundsätzen hervorgeht,
eine Zerfällung der \gesperrt{Galois}'schen Resolvente. Nehmen wir an,
es sei (wie im §. 145) $g(t) = 0$ die \gesperrt{Galois}'sche Resolvente und
$\varrho$ eine ihrer Wurzeln, und durch die Permutationen von $Q$ gehe
$\varrho$ über in

\[
\varrho, \varrho_1, \varrho_2 \ldots \varrho_{q-1},
\]

und durch die Permutationen der Nebengruppe $Q\pi_i$ in

% [sic] The source omits the comma after \varrho_{0,i} in this parallel list.
\[
\varrho_{0,i}\ \varrho_{1,i}, \varrho_{2,i} \ldots \varrho_{q-1,i},
\]

so ist jeder der Factoren von $g(t)$

\[
g_i(t) = (t - \varrho_{0,i})(t - \varrho_{1,i}) \ldots (t - \varrho_{q-1,i})
\]

durch die Substitutionen von $Q$ ungeändert und also in $\Omega(\psi)$
enthalten.

Setzen wir

\[
g(t, \psi) = (t - \varrho)(t - \varrho_1) \ldots (t - \varrho_{q-1}),
\]

so ergiebt sich, wenn wir darin eine Permutation aus $Q\pi_i$ ausführen,

\[
g_i(t) = g(t, \psi_i),
\]

und $g(t)$ ist also in folgender Art zerlegbar:

\weq{6}{g(t) = g(t, \psi)\, g(t, \psi_1) \ldots g(t, \psi_{j-1}),}

worin jeder Factor vom Grade $q$ ist, und keine zwei dieser
Factoren einen gemeinschaftlichen Theiler haben.

Sind $Q$ und $Q'$ zwei verschiedene Theiler der Gruppe $P$, so
werden $Q$ und $Q'$ gewisse Permutationen gemein haben, unter
denen sich immer die identische Permutation findet. Es ist
möglich, dass dies die einzige gemeinsame Permutation von $Q$
und $Q'$ ist, und dann heissen diese beiden Gruppen \gesperrt{theilerfremd}.
Es können aber noch mehr gemeinsame Elemente vorhanden
%CONT>>
